\documentclass[10pt,a4paper]{amsart}
\usepackage[margin=19mm]{geometry}
\usepackage{amsmath,amssymb,mathtools}
\usepackage[scr=rsfs,cal=euler]{mathalfa}
\usepackage{xfrac}
\usepackage[unicode,hidelinks,bookmarks=false]{hyperref}
\newcommand{\ie}{i.e.\ }
\newcommand{\cf}{cf.\ }
\newcommand{\resp}{respectively\ }
\newcommand{\Subsection}{\subsection}
\providecommand{\nobreakdash}{\nobreak\hbox{-}}
\setlength{\emergencystretch}{2em}
\multlinegap0pt
\usepackage{bm}
\usepackage{bbm}
\usepackage{dsfont}
\usepackage{xspace}
\usepackage{cancel}
\usepackage{mathtools}
\usepackage{amsthm}
\makeatletter
\@ifundefined{theorem}{\newtheorem{theorem}{Theorem}}{}
\@ifundefined{definition}{\newtheorem{definition}{Definition}}{}
\@ifundefined{lemma}{\newtheorem{lemma}[theorem]{Lemma}}{}
\makeatother
\newcommand\norm[1]{\left\lVert#1\right\rVert}
\newcommand{\R}{{\mathbb R}}
\title[Towards multi-purpose locally differentially-private synthet]{Towards multi-purpose locally differentially-private synthetic data release via spline wavelet plug-in estimation}
\author{Thibault Randrianarisoa, Lukas Steinberger, Botond Szab\'o}
\date{Source: arXiv:2508.13969v2 (CC BY 4.0)}
\begin{document}
\maketitle

\noindent\emph{Source and licence.} Towards multi-purpose locally differentially-private synthetic data release via spline wavelet plug-in estimation. Version arXiv:2508.13969v2 (CC BY 4.0), distributed under the Creative Commons Attribution 4.0 International licence, as recorded in the arXiv licence metadata for this exact version and shown on the abstract page https://arxiv.org/abs/2508.13969v2. Full rights notice: licenses/arxiv-2508.13969-CC-BY-4.0.txt.\par\smallskip

\noindent\textit{Editorial scope.} Counted unit: the Lemma whose source label is \texttt{lemma:W0} in the article (source0.tex lines 889--891, source label \texttt{lemma:W0}), delivered together with the article's own complete original proof of it (source0.tex lines 892--929, one \texttt{proof} environment of 38 lines). No mathematics is written, repaired or re-derived: statement and proof are the article's own text line for line. Required context retained uncounted: eq: form functional (lines 254--256); def:ND (lines 319--325). Every definition, display and label of those lines is retained unchanged. Omitted from this package and not redistributed: 1--253, 257--318, 326--888, 930--1915. Nothing inside a retained excerpt is omitted or altered: every retained body is delivered exactly as the article's lines are written, so no omission receipt of any kind accompanies this package. Citations: the retained excerpts carry no citation command at all, so no bibliography entry of the article is transcribed and no reference is redistributed. Reference adaptation: none was needed and none was made; every reference inside the delivered unit resolves inside it, so no unresolved reference or citation is produced. Printed slips kept literal: no printed word, symbol, hypothesis or number of the counted statement or of its proof was altered. Layout and class: the article's own class is \texttt{imsart} with packages including amsthm, amsmath, amsfonts, amssymb, mathtools, dsfont, footmisc, natbib, hyperref, hypernat, verbatim, textcomp, enumerate, graphicx; the wrapper is the lane's pinned \texttt{amsart} editorial wrapper at 10pt on A4, which restates the theorem-like environments the retained text needs and adds the article's own macro definitions that the retained text needs (\texttt{\textbackslash R}, \texttt{\textbackslash norm}), with extra wrapper packages (bm, bbm, dsfont, xspace, cancel, mathtools). The delivered numbering is the unit's own. Assets: no retained span contains a figure environment, a graphics-inclusion command, a PDF-page inclusion, a table or a TikZ picture; no image file is copied into this package, the package declares no asset dependency and no image file is redistributed with it. Licence: version arXiv:2508.13969v2 of this article is distributed under the Creative Commons Attribution 4.0 International licence, as recorded in the arXiv licence metadata for this exact version and shown on the abstract page https://arxiv.org/abs/2508.13969v2. A full rights notice is delivered at licenses/arxiv-2508.13969-CC-BY-4.0.txt. Characters: every retained excerpt is pure ASCII, so the delivered engine, XeLaTeX with its default Latin Modern text font, reports no missing character.
\begin{equation}\label{eq: form functional}
    \Lambda(f+h) = \Lambda(f) + T_f(h) + O(\norm{h}^2_{(2,m,\lambda)}),
\end{equation}

\begin{definition}\label{def:ND}
    Let $\mathcal W\subseteq \mathcal W_p$ and $\Lambda:\mathcal W \to \mathbb R$ be differentiable as in \eqref{eq: form functional} and atomic of index $s$. We say that $(\mathcal W, \Lambda)$ is \emph{non-degenerate} if for all $x_0\in[0,1]$ there exists a compactly supported function $\psi\in C^p(\R)$ with $\psi^{(s)}(0)\neq 0$ and a positive constant $\delta_0>0$, such that the following holds true: For all $f_0\in\mathcal W$ there exists $\eta_0>0$, such that for all $\eta\in(0,\eta_0)$ and all $\delta\in(0,\delta_0)$ the function
    $$
    f_1(x) := f_0(x) + \eta\delta^p\psi((x-x_0)/\delta), \quad x\in[0,1],
    $$
    belongs to $\mathcal W$. If this is not the case we call $(\mathcal W,\Lambda)$ \emph{degenerate}.
\end{definition}

\begin{lemma}\label{lemma:W0}
    Let $\mathcal W^\circ := \{f\in\mathcal W_p : f>0\}$ and $\Lambda:\mathcal W_p \to \mathbb R$ be differentiable as in \eqref{eq: form functional} and atomic of index $s$, $0\le s\le p$. Then $(\mathcal W^\circ,\Lambda)$ is non-degenerate as in Definition~\ref{def:ND}.
\end{lemma}

\begin{proof}
    Define $g(x):= \exp\left(-\frac{1}{1-4x^2} \right)$,  if $x\in(-\frac12,\frac12)$ and $g(x) := 0$, else, which is infinitely many times differentiable on $\R$. Notice that there must be an $x^*\in(-\frac12,\frac12)$, such that $g^{(s)}(x^*)\neq 0$, or otherwise $g$ is a polynomial on that interval. By symmetry we may assume that $x^*\in(-\frac12,0]$. For $b>0$, set $\varphi(x) := -g(x+\frac12) + bg(x-\frac12)$, which has compact support $[-1,1]$. If $x_0\in(0,1)$, we take $b=1$. If $x_0=0$, we take $b=(\int_{x^*}^{1/2}g(u)du )/(\int_{-1/2}^{1/2} g(u)du )$. Finally, if $x_0=1$, we set $b=(\int_{-1/2}^{1/2} g(u)du )/(\int_{-1/2}^{-x^*}g(u)du )$. Furthermore, if $x_0\in[0,1)$, define $\psi(x) := \varphi(x+x^*-\frac12)$ and if $x_0=1$ set it equal to $\psi(x) := \varphi(x-x^*+\frac12)$. 
    In any case, this entails that $\psi$ is compactly supported with $\psi\in C^p(\R)$ and $\psi^{(s)}(0) \neq 0$. Take $\delta_0=(x_0\land(1-x_0))/2$ if $x_0\in(0,1)$ and set $\delta_0=\frac12$ if $x_0\in\{0,1\}$. Now fix $f_0\in \mathcal W^\circ$, which entails that $\inf_{x\in[0,1]}f_0(x)>0$, and $\|f_0\|_{(\infty,p)}<M$, and define $\eta_0 = \left(\frac{\inf_{x\in[0,1]}f_0(x)}{\|\psi\|_\infty} \right)\land\left( \frac{M-\|f_0\|_{(\infty,p)}}{\|\psi\|_{(\infty,p)}}\right)>0$. For $0<\eta<\eta_0$, $0<\delta<\delta_0$ and $x\in[0,1]$ consider 
    $$
    f_1(x) := f_0(x) + \eta\delta^p\psi((x-x_0)/\delta).
    $$
    By construction, we have $\inf_{x\in[0,1]}f_1(x)\ge \inf_{x\in[0,1]}f_0(x) - \eta\|\psi\|_{L^\infty}>0$ and $\|f_1\|_{(\infty,p)} \le \|f_0\|_{(\infty,p)} + \eta\|\psi\|_{(\infty,p)} < M$. Clearly $f_1\in C^p[0,1]$. In the case $x_0\in[0,1)$, notice that
    \begin{align*}
        \int_0^1 \psi\left( \frac{x-x_0}{\delta}\right)dx
        &=\delta \int_{-\frac{x_0}{\delta}}^{\frac{1-x_0}{\delta}} \psi(u)du
        =\delta
        \int_{-\frac{x_0}{\delta}+x^*-\frac12}^{\frac{1-x_0}{\delta}+x^*-\frac12} \varphi(u)du\\
        &=
        -\delta \int_{-\frac{x_0}{\delta}+x^*}^{\frac{1-x_0}{\delta}+x^*}g(u)du + \delta b \int_{-\frac{x_0}{\delta}+x^*-1}^{\frac{1-x_0}{\delta}+x^*-1}g(u)du.
    \end{align*}
    If we specialize to the case $x_0\in(0,1)$, we observe that the integral limits contain the support of $g$ whenever $\delta<(x_0\land(1-x_0))/2$, and since $b=1$ in that case, we see that the expression in the previous display is equal to zero. In the case $x_0=0$ we see that for $\delta<\frac12$ the expression in the previous display reads
    $$
    -\delta \int_{x^*}^{\frac{1}{\delta}+x^*}g(u)du + \delta b \int_{x^*-1}^{\frac{1}{\delta}+x^*-1}g(u)du
    =
    -\delta \int_{x^*}^{\frac12}g(u)du + \delta b \int_{-\frac12}^{\frac12}g(u)du = 0,
    $$
    in view of our definition of $b$.
    A similar consideration for the case $x_0=1$ reveals that
    \begin{align*}
        \int_0^1 \psi\left( \frac{x-x_0}{\delta}\right)dx
        &=\delta \int_{-\frac{1}{\delta}}^{0} \psi(u)du
        =\delta
        \int_{-\frac{1}{\delta}-x^*+\frac12}^{-x^*+\frac12} \varphi(u)du\\
        &=
        -\delta \int_{-\frac{1}{\delta}-x^*+1}^{-x^*+1}g(u)du + \delta b \int_{-\frac{1}{\delta}-x^*}^{-x^*}g(u)du.
    \end{align*}
    Thus, if $\delta<\frac12$, this equals 
    $$
    -\delta \int_{-\frac12}^{\frac12}g(u)du + \delta b \int_{-\frac12}^{-x^*}g(u)du = 0
    $$
    in view of our choice of $b$ in that case.
    Hence, in all the cases, we have shown that $f_1\in\mathcal W^\circ$.
\end{proof}

\end{document}
